\section{Item d)}

A interpretação do espectro harmônico extraído do FEMM deve ser feita com cuidado porque o eixo horizontal dos gráficos representa a ordem harmônica espacial mecânica, denotada por $n$, e não diretamente a ordem harmônica elétrica do enrolamento. Como a máquina analisada possui 6 polos, há $p=3$ pares de polos. Dessa forma, a relação entre a harmônica elétrica $\nu$ da distribuição de enrolamento e a ordem mecânica observada na FFT é dada por

$$n = \nu p.$$

Consequentemente, a componente fundamental da onda de densidade de fluxo no entreferro, correspondente à harmônica elétrica $\nu=1$, aparece na ordem mecânica

$$n = 1 \cdot 3 = 3.$$

Isso é exatamente o que se observa no espectro obtido: a maior componente está em $n=3$, com amplitudes próximas de $0{,}79\,\mathrm{T}$ para as três configurações. Os valores extraídos foram aproximadamente $0{,}7943\,\mathrm{T}$ para camada simples com passo pleno, $0{,}7930\,\mathrm{T}$ para camada dupla com passo pleno e $0{,}7764\,\mathrm{T}$ para camada dupla com passo encurtado. Portanto, o resultado confirma que o processamento está identificando corretamente a fundamental espacial de uma máquina de 6 polos.

A principal fonte de erro de interpretação ocorre quando se observa diretamente as barras $n=5$ e $n=7$ do gráfico e se conclui que elas representam a 5ª e a 7ª harmônicas elétricas. Essa interpretação não é correta. Como o eixo está em ordem mecânica, as barras $n=5$ e $n=7$ não correspondem às harmônicas elétricas $\nu=5$ e $\nu=7$. Para uma máquina com $p=3$, essas posições corresponderiam formalmente a $\nu = 5/3$ e $\nu = 7/3$, que não são as harmônicas inteiras usualmente analisadas na teoria de enrolamentos trifásicos. Assim, os valores observados em $n=5$ e $n=7$ não devem ser usados para validar a redução esperada dos harmônicos de 5ª e 7ª ordem do enrolamento.

As harmônicas elétricas de interesse devem ser procuradas nas ordens mecânicas obtidas pela conversão $n=\nu p$. Logo, para a 5ª harmônica elétrica,

$$n = 5 \cdot 3 = 15,$$

e, para a 7ª harmônica elétrica,

$$n = 7 \cdot 3 = 21.$$

Portanto, a validação do efeito do passo encurtado deve ser feita nas barras $n=15$ e $n=21$, e não nas barras $n=5$ e $n=7$.

Os resultados obtidos confirmam a tendência teórica esperada. Para a ordem mecânica $n=15$, correspondente à 5ª harmônica elétrica, as amplitudes foram aproximadamente $0{,}0200\,\mathrm{T}$ na camada simples com passo pleno, $0{,}0201\,\mathrm{T}$ na camada dupla com passo pleno e apenas $0{,}00269\,\mathrm{T}$ na camada dupla com passo encurtado. Em termos relativos à fundamental de cada caso, isso corresponde a cerca de $2{,}52\%$, $2{,}53\%$ e $0{,}35\%$, respectivamente. Assim, ao passar de camada dupla com passo pleno para camada dupla com passo encurtado, a 5ª harmônica elétrica foi reduzida para aproximadamente $13{,}7\%$ do valor relativo original, isto é, uma redução de cerca de $86\%$.

De modo semelhante, para a ordem mecânica $n=21$, correspondente à 7ª harmônica elétrica, as amplitudes foram aproximadamente $0{,}0185\,\mathrm{T}$ na camada simples com passo pleno, $0{,}0185\,\mathrm{T}$ na camada dupla com passo pleno e $0{,}00314\,\mathrm{T}$ na camada dupla com passo encurtado. Em termos relativos à fundamental, esses valores correspondem a cerca de $2{,}33\%$, $2{,}33\%$ e $0{,}40\%$, respectivamente. Portanto, a 7ª harmônica elétrica foi reduzida para aproximadamente $17{,}3\%$ do valor relativo observado no caso de camada dupla com passo pleno, o que representa uma redução de cerca de $83\%$.

Esses resultados estão qualitativamente de acordo com a teoria de enrolamentos. Para um passo encurtado de $5/6$ do passo polar, o fator de passo reduz fortemente as componentes de 5ª e 7ª ordem. Usando a formulação clássica,

$$k_{p,\nu} = \cos\left(\frac{\nu \varepsilon}{2}\right),$$

com encurtamento elétrico $\varepsilon = 30^\circ$, obtêm-se

$$|k_{p,5}| = |\cos(75^\circ)| \approx 0{,}2588,$$

e

$$|k_{p,7}| = |\cos(105^\circ)| \approx 0{,}2588.$$

Ou seja, a teoria prevê uma forte atenuação das harmônicas elétricas de 5ª e 7ª ordem quando se utiliza passo encurtado. O FEMM reproduziu exatamente essa tendência, pois as componentes em $n=15$ e $n=21$ diminuíram de forma expressiva na configuração de camada dupla com passo encurtado.

A diferença entre a redução teórica ideal e a redução observada numericamente é esperada. A teoria dos fatores de enrolamento considera, em geral, uma distribuição idealizada de força magnetomotriz, entreferro uniforme, permeabilidade elevada ou infinita no ferro e ausência de efeitos de ranhura, saturação e rotor. Já o modelo FEMM considera a geometria ranhurada, a permeabilidade finita e possivelmente não linear do aço, a presença do rotor e a extração da componente radial $B_r$ em um raio fixo do entreferro. Esses fatores modificam as amplitudes finais dos harmônicos de densidade de fluxo, especialmente quando se comparam amplitudes pequenas.

Além disso, o espectro apresenta componentes importantes em ordens associadas às ranhuras. Considerando $S=36$ ranhuras de estator e $p=3$, as ordens $S-p=33$ e $S+p=39$ são candidatas naturais a harmônicos de ranhura. De fato, as amplitudes nessas ordens são significativamente maiores que as dos harmônicos de baixa ordem do enrolamento, especialmente em $n=33$. Isso indica que parte importante da ondulação espacial de $B_r$ no entreferro não vem apenas da distribuição senoidal ideal do enrolamento, mas também da modulação de permeância causada pelo ranhuramento.

Portanto, a interpretação correta dos resultados é que o passo encurtado funcionou conforme o esperado. A conclusão equivocada surgiria apenas se as barras $n=5$ e $n=7$ fossem interpretadas como 5ª e 7ª harmônicas elétricas. Como a FFT está em ordem mecânica, a comparação correta deve ser feita em $n=15$ e $n=21$. Nessas ordens, os resultados FEMM mostram redução acentuada das componentes correspondentes às harmônicas elétricas de 5ª e 7ª ordem, validando qualitativamente o comportamento previsto pela teoria dos fatores de enrolamento.
